\chapter{Thật và Giả}
\markboth{Thật và Giả}{Real and False}

\section{Một người đăng rất vui nhưng sống rất mệt.}
\begin{truyen}{Một người đăng rất vui nhưng sống rất mệt.}{A bright image can hide a tired life.}
\chuhoa{M}{ạng xã hội của cô đầy ảnh cười, câu tích cực và những chuyến đi đẹp. Nhưng ngoài đời cô mất ngủ, lo âu, và không biết chia sẻ thật với ai. Cái giả không phải lúc nào cũng là nói dối thô; nhiều khi nó là sống quá lâu trong một hình ảnh bóng bẩy.}
\textit{(Her social media was full of smiling photos, positive lines, and beautiful trips. In real life she had insomnia, anxiety, and no one with whom she could be honest. Falsehood is not always crude lying; often it is living too long inside a polished image.)}
\end{truyen}
\begin{baihoc}
Sống giả lâu ngày không chỉ làm người khác hiểu nhầm, mà còn làm chính mình lạc khỏi cảm xúc thật.
\textit{(Living falsely for too long not only misleads others, but also separates us from our own real feelings.)}
\end{baihoc}
\begin{ghinhoanh}
\textbf{Short English to remember:}

A polished image can hide pain.
\end{ghinhoanh}
\begin{tuhocnhanh}
1. Em có đang duy trì một hình ảnh khác xa thực tế của mình không? \textit{(Are you maintaining an image that is far from your real life?)}

2. Em có một nơi nào để được sống thật không? \textit{(Do you have any place where you can live honestly?)}
\end{tuhocnhanh}
\ngancach

\section{Một lời xin lỗi thật hơn nhiều món quà giả vờ.}
\begin{truyen}{Một lời xin lỗi thật hơn nhiều món quà giả vờ.}{A true apology weighs more than decorated gestures.}
\chuhoa{S}{au khi làm bạn buồn, anh mua quà nhưng vẫn né nói lời nhận lỗi. Người bạn nhận quà mà lòng không nhẹ hơn. Đến khi anh thật sự nói ra: “Mình sai ở chỗ này”, mối quan hệ mới có cơ hội lành lại.}
\textit{(After hurting a friend, he bought gifts but avoided admitting the wrong. The friend received the presents without feeling lighter. Only when he truly said, “I was wrong here,” did the relationship have a chance to heal.)}
\end{truyen}
\begin{baihoc}
Điều thật thường đơn sơ nhưng có sức chữa lành mạnh hơn vẻ đẹp bề ngoài.
\textit{(What is real is often simple, yet more healing than outer decoration.)}
\end{baihoc}
\begin{ghinhoanh}
\textbf{Short English to remember:}

Truth heals more than performance.
\end{ghinhoanh}
\begin{tuhocnhanh}
1. Em có chuyện gì nên nói thật thay vì vòng vo không? \textit{(Is there something you should say honestly instead of circling around?)}

2. Em có đang dùng hành động đẹp để né một sự thật khó nói? \textit{(Are you using nice gestures to avoid a hard truth?)}
\end{tuhocnhanh}
\ngancach

\section{Học thật chậm nhưng bền hơn học giả nhanh.}
\begin{truyen}{Học thật chậm nhưng bền hơn học giả nhanh.}{Real learning may be slower, but it stays.}
\chuhoa{N}{gười bạn kia thuộc bài rất nhanh trước giờ kiểm tra, nhưng vài hôm sau quên sạch. Còn em học chậm hơn, phải hỏi nhiều hơn, nhưng hiểu được tận gốc. Cái giả thường cho cảm giác nhanh thắng, còn cái thật lặng lẽ đi xa.}
\textit{(One student memorized quickly before tests but forgot everything a few days later. Another learned more slowly, asked more questions, yet understood the roots. The false often feels like a quick win, while the real travels farther in silence.)}
\end{truyen}
\begin{baihoc}
Học thật không phải lúc nào cũng cho kết quả nhanh, nhưng nó cho kết quả còn ở lại.
\textit{(Real learning does not always produce quick results, but it produces results that remain.)}
\end{baihoc}
\begin{ghinhoanh}
\textbf{Short English to remember:}

Real learning stays longer.
\end{ghinhoanh}
\begin{tuhocnhanh}
1. Em đang học để thật hiểu hay chỉ để qua trước mắt? \textit{(Are you learning to truly understand or only to pass for the moment?)}

2. Kiến thức nào của em hiện đang chỉ là vẻ ngoài? \textit{(Which part of your knowledge is still mostly surface?)}
\end{tuhocnhanh}
\ngancach

\section{Giả mạnh để che một nỗi sợ thật.}
\begin{truyen}{Giả mạnh để che một nỗi sợ thật.}{False strength often hides true fear.}
\chuhoa{C}{ó người nói rất cứng, luôn tỏ ra bất cần, nhưng chỉ cần bị chạm vào lòng tự trọng là phản ứng dữ dội. Đằng sau vẻ bất cần ấy không phải lúc nào cũng là bản lĩnh; nhiều khi chỉ là một nỗi sợ bị xem thường chưa được chữa lành.}
\textit{(Some people speak harshly and act as though nothing matters, yet react violently whenever their pride is touched. Behind that indifference there is not always courage; often there is an unhealed fear of being looked down upon.)}
\end{truyen}
\begin{baihoc}
Muốn sống thật, phải dám nhận cả phần mong manh của mình.
\textit{(To live truthfully, we must dare to admit our fragile parts.)}
\end{baihoc}
\begin{ghinhoanh}
\textbf{Short English to remember:}

False strength hides fear.
\end{ghinhoanh}
\begin{tuhocnhanh}
1. Em có đang dùng một vẻ ngoài cứng cỏi để che điều gì bên trong không? \textit{(Are you using a tough exterior to hide something within?)}

2. Nếu sống thật hơn, em sẽ phải đối diện với nỗi sợ nào? \textit{(If you lived more honestly, which fear would you need to face?)}
\end{tuhocnhanh}
\ngancach

\section{Điều thật thường ít hào nhoáng nhưng đáng tin hơn.}
\begin{truyen}{Điều thật thường ít hào nhoáng nhưng đáng tin hơn.}{What is real rarely needs loud decoration.}
\chuhoa{M}{ột người thầy không nói câu nào quá kêu, cũng không tự tô mình thành hình mẫu. Nhưng thầy giữ lời, dạy kỹ, đối xử công bằng, và có mặt khi học sinh cần. Sự thật nhiều khi không chói lòa; nó bền bỉ, giản dị và khiến người ta tin.}
\textit{(One teacher never said anything overly grand and never painted himself as a perfect model. Yet he kept promises, taught carefully, treated students fairly, and showed up when needed. Reality is often not dazzling; it is durable, simple, and trustworthy.)}
\end{truyen}
\begin{baihoc}
Đáng tin quan trọng hơn ấn tượng.
\textit{(Being trustworthy matters more than being impressive.)}
\end{baihoc}
\begin{ghinhoanh}
\textbf{Short English to remember:}

Trust beats glitter.
\end{ghinhoanh}
\begin{tuhocnhanh}
1. Em thường bị thu hút bởi sự hào nhoáng hay sự đáng tin? \textit{(Are you more attracted to glitter or to trustworthiness?)}

2. Em muốn người khác nhớ mình vì vẻ ngoài hay vì độ thật? \textit{(Do you want others to remember you for appearance or for authenticity?)}
\end{tuhocnhanh}
\ngancach
